\paragraph{Theorem F.2.31 (Proved: complete fundamental-group data for every input).}
Let $k=2$ or $3$, and let an integral rank-six tuple satisfy the
finite-order and primitive maximal-cusp conditions of F.2.  No
irreducibility, determinant-one, or prospective-survivor hypothesis is
imposed.  Specify each actual finite affine action as
\[
 \tau_i(x)=M_ix+v_i/m_i,\qquad v_i\in\Lambda=\mathbb Z^6,
 \quad \tau_i^{m_i}=1\text{ on }\Lambda_{\mathbb R}/\Lambda,
 \tag{F.G1}
\]
and require its cyclic action to be free.  The numerator need not be
fixed by $M_i$.  Define its actual integral power vector and the cusp
lattice by
\[
 w_i=\frac1{m_i}\sum_{j=0}^{m_i-1}M_i^jv_i\in\Lambda^{M_i},
 \qquad U=\operatorname{im}\bigl((M_1\cdots M_k)^{-1}-I\bigr).
 \tag{F.G2}
\]
Use left conjugation transport for the peripheral generators $s_i$,
and record the integral translation clutching by $\kappa\in\Lambda$
in $s_1\cdots s_kx^\kappa=1$.  Then the full fundamental group is
\[
 \left\langle x^\lambda,s_1,\ldots,s_k\ \left|\
 \begin{array}{l}
 x^\lambda x^\mu=x^{\lambda+\mu},\quad x^u=1\ (u\in U),\\
 s_ix^\lambda s_i^{-1}=x^{M_i\lambda},\quad
 s_i^{m_i}=x^{w_i},\quad s_1\cdots s_kx^\kappa=1
 \end{array}\right.\right\rangle .                         \tag{F.G3}
\]
Only a basis of $\Lambda$ and of $U$ is needed in this finite
presentation.  The program below also determines exactly which
quotient of the apparent fibre subgroup injects.

Put
\[
 L=\sum_{g\in\langle M_1,\ldots,M_k\rangle}gU,
 \qquad A=\Lambda/L.                                      \tag{F.G4}
\]
This is an integral sum, without saturation.  For two points, put
$d=\gcd(m_1,m_2)$, $M=M_1$, and
$S_j(M)=I+M+\cdots+M^{j-1}$.  Choose integers $a,b$ with
$am_1+bm_2=d$, and set, in $A$,
\[
 \begin{split}
 w'_2&=-w_2-S_{m_2}(M)\kappa,\\
 e&=(m_2/d)w_1-(m_1/d)w'_2,\qquad
 z=aw_1+bw'_2,\qquad \overline A=A/\langle e\rangle .
 \end{split}                                             \tag{F.G5}
\]
There is an exact sequence and a complete marked extension description
\[
 1\longrightarrow\overline A\longrightarrow\pi_1
   \longrightarrow C_d\longrightarrow1,
 \qquad txt^{-1}=M(x),\quad t^d=z.                         \tag{F.G6}
\]
Thus its order is $d|\overline A|$ if $\overline A$ is finite;
otherwise it is infinite.  In particular $d=1$ gives precisely
the abelian group $\overline A$.

For three genuine points, so $m_i\ge2$, let
\[
 \Gamma=\langle q_1,q_2,q_3\mid q_i^{m_i}=1,
                         \ q_1q_2q_3=1\rangle,
 \qquad \epsilon=\sum_i1/m_i-1.                           \tag{F.G7}
\]
The action $q_i\mapsto M_i$ is well defined on $A$.
If $\epsilon\le0$, the exact sequence is
\[
                    1\longrightarrow A\longrightarrow\pi_1
                         \longrightarrow\Gamma\longrightarrow1. 
 \tag{F.G8}
\]
If $\epsilon>0$, $\Gamma$ is the finite spherical triangle group.
Using its action on $A$, form the exact transfer vector
\[
 e=\sum_{g\in\Gamma}g\kappa+
       \sum_{i=1}^3\ \sum_{g\in\Gamma/\langle q_i\rangle}g w_i.
 \tag{F.G9}
\]
The second sums use left coset representatives and are independent
of their choice.  In this case the exact sequence has kernel
$A/\langle e\rangle$, with the same lifted relations (F.G3).
Consequently (F.G3), (F.G6), (F.G8), and (F.G9) give every requested
fundamental group by finite operations, including the precise
injected abelian kernel and all extension data.  These statements
do not claim a conjugacy classification of the original tuples.

For every input, the full first integral homology is the Smith
cokernel of the integer matrix
\[
 \begin{pmatrix}
 M_1-I&\cdots&M_k-I&-w_1&\cdots&-w_k&\kappa\\
 0&\cdots&0&&\operatorname{diag}(m_1,\ldots,m_k)&&\boldsymbol1
 \end{pmatrix}.                                           \tag{F.G10}
\]
Here the bottom left zero block has $k$ rows.  The code records
the full matrix, not only its invariant factors.

\begin{proof}
\emph{Actual finite and cusp caps.}
The affine lift on $\mathbb R^6$ has $m_i$th power translation by
$w_i$, as the displayed sum shows.  It commutes with that power,
so $M_iw_i=w_i$.  The group generated by lattice translations
and this lift is the fundamental group of the reduced flat
six-manifold.  Its presentation consists of the lattice, its
$s_i$ conjugation relations, and $s_i^{m_i}=w_i$.
The finite cap deformation retracts onto that reduced manifold
by the equivariant radial contraction of its root disc.
Thus these are the actual cap relations, even for a numerator
which is not an invariant lattice vector.

For the primitive Kuhn cap, the compact toric strata in its
normalization and all their intersection components are simply
connected.  Applying van Kampen successively to their tubular
neighborhoods identifies the fundamental group of the central
fibre with that of its dual three-torus.  The nearby fibre map
kills the angular lattice $U$ and sends the long-period lattice
$\Lambda/U$ isomorphically onto this dual-torus group.  This is
also the explicit toric collapse used in E.1.3.  The primitive
index-one condition gives $U=\ker(M_0-I)$ and makes $\Lambda/U$
free of rank three.  A section through a smooth central stratum
bounds the base meridian by its parameter disc.  The smooth cap
retracts onto the central fibre, so it adds exactly these
relations, with no extra meridian relation left unspecified.

Over the punctured sphere the bundle has a section up to homotopy,
because the base retracts to a graph.  Its group is the appropriate
semidirect product.  Choose the cap sections and retain their
relative translation clutching as $\kappa$.  Van Kampen with the
actual boundary maps just described gives exactly (F.G3).
Changing the transport convention causes no group ambiguity:
reverse the ordered list and replace each affine map by its
inverse, and replace $\kappa$ by $-\kappa$.  The new linear
matrices are $M_k^{-1},\ldots,M_1^{-1}$, the power vectors are
$-w_k,\ldots,-w_1$, and the cusp image lattice is unchanged.
The substitutions $s_i\mapsto s_{k+1-i}^{-1}$ give the explicit
isomorphism between these two presentations.  This supplies the
dictionary when geometric transport is recorded with opposite
path-composition order.

\emph{Integral closure and two-point elimination.}
In a normal abelian fibre subgroup, the normal closure of $U$
is precisely the integer sum (F.G4).  Start with its column
Hermite basis and repeatedly append its images under every
$M_i$ and $M_i^{-1}$, then take Hermite normal form again.
If the rational rank increases, it can do so at most three
times.  At the first nonincrease the rational span is invariant
under every generator and its inverse, hence never increases
again.  If the current lattice has saturation index $D$ in
that span, every subsequent strict enlargement divides this
index by an integer at least two.  There can therefore be at
most $\lfloor\log_2D\rfloor$ such enlargements.  This proves
termination and gives the computable bound used by the code;
it never replaces a finite-index lattice by its saturation.

On $A$ one has $M_2=M^{-1}$ in the two-point case.  The product
relation eliminates $s_2=t^{-1}x^{-\kappa}$, where $t=s_1$.
Multiplying this word $m_2$ times gives
\[
        (t^{-1}x^{-\kappa})^{m_2}
            =t^{-m_2}x^{-S_{m_2}(M)\kappa}.
\]
Thus the two power relations become $t^{m_1}=w_1$ and
$t^{m_2}=w'_2$.  Both right sides are fixed by $M$ on $A$.
Also $M^{m_1}=M^{m_2}=1$, hence $M^d=1$ there.
Bezout elimination gives exactly $e=0$ and $t^d=z$ in
(F.G5).  Replacing the Bezout pair changes $z$ by a multiple
of $e$ and therefore changes no extension datum.

To prove the asserted injection, construct
$\overline A\rtimes_M\mathbb Z$ and quotient by its central
element $t^dz^{-1}$.  This cyclic subgroup has nonzero
projection $d\mathbb Z$ to the infinite cyclic factor, so it
intersects $\overline A$ trivially.  The quotient has precisely
the presentation obtained by elimination.  This proves (F.G6),
including the absence of any further hidden fibre relation.

\emph{The three-point kernel and the finite transfer.}
For completeness this step computes the fibre injection; it
does not infer it from the appearance of a normal subgroup in
a presentation.  Let $\mathcal B$ be the orbifold classifying
space obtained from the three-holed sphere by attaching, at
its three boundary circles, the mapping cylinders of
$S^1\to BC_{m_i}$.  Over the punctured sphere take the bundle
with fibre $BA=K(A,1)$ and the specified action.  For each
boundary, the cyclic extension
\[
 E_i=\langle A,s_i\mid s_i a s_i^{-1}=M_i(a),\ s_i^{m_i}=w_i\rangle
\]
contains $A$ injectively: the same semidirect-product construction
as above applies because $M_i^{m_i}=1$ and $M_iw_i=w_i$.
The associated fibration $BA\to BE_i\to BC_{m_i}$ restricts
on $S^1$ to the required mapping-torus fibration.  Replace the
maps by their mapping-path fibrations before homotopy gluing,
so that the fibre is preserved on every attaching cylinder.
Glue these three fibrations, using the specified relative clutching
$\kappa$ on the punctured sphere.  Van Kampen gives (F.G3)
for the resulting total space.  Thus the homotopy sequence
of this fibration over $\mathcal B$ calculates precisely
the image of $A$ in the group under discussion.

When $\epsilon\le0$, the universal orbifold cover is the plane
or hyperbolic plane.  Its contractibility identifies
$\mathcal B$ with $B\Gamma$: indeed its Borel construction
has the just-described boundary-cylinder model and contractible
universal fibre.  The Euclidean and hyperbolic alternatives
are the oriented two-orbifold classification in
\cite[Theorems 2.3--2.4]{rootScott83}.
Since $\pi_2(B\Gamma)=0$, the homotopy sequence proves (F.G8).

When $\epsilon>0$, the universal orbifold cover is $S^2$.
Pulling the constructed $BA$ fibration to this finite cover
trivializes its local system.  Its transgression
$\pi_2(S^2)\to A$ is its integral clutching vector.
There is one lifted ordinary base cell for each $g\in\Gamma$,
contributing $g\kappa$, and one lifted cone cap for each left
coset of $\langle q_i\rangle$, contributing $gw_i$.
To check the sign, the base relation has fibre term $+\kappa$
and each cone relation has term $-w_i$; the oriented sphere
chain subtracts the cone cells from the base cell.  The result
is therefore the plus sign in both sums of (F.G9).
Changing a coset representative by $q_i^j$ does not change
its contribution because $w_i$ is fixed.  The result is
$\Gamma$ invariant, as also checked by the code.  The long
homotopy sequence now says that the only additional kernel
is exactly $\langle e\rangle$.  This proves the spherical
assertion, with the orientation fixed by this cell calculation.

There is no implicit infinite search for a spherical group.
After sorting its orders the possibilities are
$(2,2,n),(2,3,3),(2,3,4),(2,3,5)$, with groups of orders
$2n,12,24,60$, respectively.  These follow directly from
the positive spherical-triangle angle condition.
The code constructs the dihedral, tetrahedral, octahedral,
or icosahedral permutation group and finds the marked
peripheral generators in that finite group.  Breadth-first
words supply their actions on $A$; every generator edge is
checked modulo the actual lattice $L$.  It then computes
the two finite sums in (F.G9) and their integer Smith form.

An order-one site is not a genuine cone.  Its matrix is the
identity and its relation is $s_i=x^{w_i}$.  Eliminate it
from (F.G3); moving this vector to the right of the remaining
peripheral product replaces the clutching by
\[
       \kappa+(M_{i+1}\cdots M_k)^{-1}w_i.                 \tag{F.G11}
\]
This reduces the three-point case with such a site to the
two-point computation.  A primitive rank-three cusp cannot
remain with only one finite linear monodromy: a finite-order
unipotent matrix is the identity.  Thus no smaller case has
been omitted within the hypotheses of this theorem.

\emph{Freeness, abelianization, and executed certificates.}
Freeness of (F.G1) is itself decidable exactly.  A nontrivial
stabilizer in a cyclic group contains a subgroup of prime
order, so it is enough to check $\tau_i^{m_i/p}$ for every
prime $p\mid m_i$.  For an affine map $Rx+b$ of a torus,
a fixed point exists precisely when every integral covector
in $\ker(R^t-I)$ takes an integral value on $b$.
Smith form proves this criterion: the image of $R-I$ on
the torus is the connected subtorus annihilated by those
primitive covectors.  The code computes the powers in the
rational affine matrix group and performs these finite tests.

Abelianizing (F.G3) yields (F.G10).  No cusp column was lost:
\[
 \operatorname{im}(M_0-I)
       \subset\sum_i\operatorname{im}(M_i-I)
\]
follows by telescoping the product on the right; inversion
preserves the image of its difference from the identity.
The Smith form therefore computes every free and torsion
summand of $H_1$ exactly.

The program \texttt{fourfolds\_vankampen.py} saves all input
matrices, power vectors, integral closure traces, group
extension data, and homology presentation matrices.
It checks the nine explicitly listed three-point parameters
of its original test set, four two-point order/clutching
choices, a genuinely nonfixed affine numerator, and a
finite-index closure step.  It also checks a nonabelian
rank-six two-point tuple, a spherical $(2,4,2)$ example
with injected kernel $\mathbb Z^2\oplus\mathbb Z/4$, and
the infinite perfect hyperbolic example with orders
$(5,4,19)$.  The nonabelian test deliberately belongs to
the full $\operatorname{GL}_6$ scope and has orientation-
reversing finite matrices; it is not claimed to admit a
complex compactification.  All bounds above are derived
from each input, rather than imposed as a search cutoff.
\end{proof}
